\chapter{Symbols}
\label{app:notation}


Notation is organized by category: categories and functors, operads
and algebras, configuration spaces and forms, and chiral structures.

\begin{remark}[Two-axis \texorpdfstring{$E$}{E}-notation; Remark~\ref{rem:En-terminology}]
Bare $\En$ denotes the little-disks/topological factorization axis on
real $n$-manifolds.  The curve-based hierarchy is written only with the
suffix ``-chiral'': $\Einf$-chiral, $\Pinf$-chiral, $\Eone$-chiral.
A complex curve enters the topological ladder at $n = 2$.
\end{remark}

\section{Categories and functors}
\label{sec:notation-categories}

\begin{longtable}{p{3cm} p{10cm}}
\caption{Categories and functors} \\
\toprule
\textbf{Symbol} & \textbf{Meaning} \\
\midrule
\endfirsthead
\toprule
\textbf{Symbol} & \textbf{Meaning} \\
\midrule
\endhead
\midrule
\multicolumn{2}{r}{\textit{Continued on next page}} \\
\endfoot
\bottomrule
\endlastfoot

$\Cat$ & Category of small categories \\
$\Catinf$ & $\infty$-category of $\infty$-categories \\
$\Spc$ & $\infty$-category of spaces (Kan complexes) \\
$\Ch(k)$ & Category of chain complexes over $k$ \\
$\dgVect$ & Category of dg-vector spaces \\
$\grVect$ & Category of graded vector spaces \\
$\DMod(X)$ & $\infty$-category of D-modules on $X$ \\
$\QCoh(X)$ & $\infty$-category of quasi-coherent sheaves \\
$\IndCoh(X)$ & $\infty$-category of ind-coherent sheaves \\
$\Shv(X)$ & $\infty$-category of sheaves on $X$ \\
$\Perv(X)$ & Category of perverse sheaves \\
$\PrL$ & $\infty$-category of presentable $\infty$-categories with left adjoints \\
$\PrR$ & $\infty$-category of presentable $\infty$-categories with right adjoints \\
$\Hom(-, -)$ & Hom-set or Hom-space \\
$\RHom(-, -)$ & Derived/internal Hom \\
$\Map(-, -)$ & Mapping space in an $\infty$-category \\
$\colim$ & Colimit \\
$\holim$ & Homotopy limit \\
$\hocolim$ & Homotopy colimit \\
$f_*$, $f^*$ & Pushforward and pullback functors \\
$f_!$, $f^!$ & Exceptional pushforward and pullback \\
$\VD$ & Verdier duality functor \\
$\DR$ & De Rham functor \\
$\Sol$ & Solutions functor \\
\end{longtable}


\section{Operads and algebras}
\label{sec:notation-operads}

\begin{longtable}{p{3cm} p{10cm}}
\caption{Operads and algebras} \\
\toprule
\textbf{Symbol} & \textbf{Meaning} \\
\midrule
\endfirsthead
\toprule
\textbf{Symbol} & \textbf{Meaning} \\
\midrule
\endhead
\midrule
\multicolumn{2}{r}{\textit{Continued on next page}} \\
\endfoot
\bottomrule
\endlastfoot

$\Op$ & Category of operads \\
$\coOp$ & Category of cooperads \\
$\Opinfty$ & $\infty$-category of $\infty$-operads \\
$\Alg_\cP(\cC)$ & $\cP$-algebras in $\cC$ \\
$\CoAlg_{\mathcal{Q}}(\cC)$ & $\mathcal{Q}$-coalgebras in $\cC$ \\
$\Ass$ & Associative operad \\
$\Com$ & Commutative operad \\
$\Lie$ & Lie operad \\
$\Pois$ & Poisson operad \\
$\BV$ & Batalin--Vilkovisky operad \\
$\Grav$ & Gravity operad \\
%
$\Eone$ & $\Eone$-operad (little $1$-disks; classical/topological associative axis) \\
$\Etwo$ & $\Etwo$-operad (little $2$-disks; oriented-surface/topological axis) \\
$\En$ & Little $n$-disks operad on real $n$-manifolds \\
$\Einf$ & $\Einf$-operad (little $\infty$-disks; topological commutative axis) \\
$\Ainf$ & $\Ainf$-operad (homotopy associative) \\
%
$\Linf$ & $\Linf$-operad (homotopy Lie) \\
$\Cinf$ & $\Cinf$-operad (homotopy commutative) \\
$\Pinf$ & $\Pinf$-operad (homotopy Poisson) \\
$\cP^!$ & Linear dual operad \\
$\cP^{\scriptstyle \text{\textexclamdown}}$ & Koszul dual cooperad \\
%
$\Free(\cM)$ & Free algebra on module $\cM$ \\
$\Cofree(\cN)$ & Cofree coalgebra on comodule $\cN$ \\
$\B(\cA)$ & Bar construction of algebra $\cA$ (a dg coalgebra; see $\Bbar$ for the reduced version) \\
$\Cobar(\cC)$ & Cobar construction of coalgebra $\cC$ (a dg algebra; $\Cobar(\Bbar(\cA)) \xrightarrow{\sim} \cA$ when $\cA$ is Koszul) \\
$\Bbar(\cA)$ & Reduced bar construction $= (T^c(s^{-1}\bar{\cA}),\, d)$: a conilpotent dg coalgebra.  This is the Koszul dual \emph{coalgebra} $\cA^{\text{\textexclamdown}}$; dualizing gives the Koszul dual \emph{algebra} $\cA^! = \Bbar(\cA)^\vee$ \\
%
$\overline{\Cobar}(\cC)$ & Reduced cobar construction \\
$\Bbar^{\mathrm{geom}}$ & Geometric bar complex \\
$\Cobar^{\mathrm{geom}}$ & Geometric cobar complex \\
$\Tw(\cC, \cP)$ & Twisting morphisms from $\cC$ to $\cP$ \\
$\MC(\fg)$ & Maurer--Cartan elements in $\fg$ \\
$\CE(\fg)$ & Chevalley--Eilenberg complex \\
\end{longtable}


\section{Configuration spaces and forms}
\label{sec:notation-config}

\begin{longtable}{p{3cm} p{10cm}}
\caption{Configuration spaces and forms} \\
\toprule
\textbf{Symbol} & \textbf{Meaning} \\
\midrule
\endfirsthead
\toprule
\textbf{Symbol} & \textbf{Meaning} \\
\midrule
\endhead
\midrule
\multicolumn{2}{r}{\textit{Continued on next page}} \\
\endfoot
\bottomrule
\endlastfoot

$C_n(X)$ & Ordered configuration space of $n$ distinct points in $X$ \\
$B_n(X)$ & Unordered configuration space $C_n(X)/\Sigma_n$ (not to be confused with the braid group $B_n$) \\
$\overline{C}_n(X)$ & Fulton--MacPherson compactification of $C_n(X)$ \\
$\FM_n(X)$ & Alternative notation for $\overline{C}_n(X)$ \\
$\Ran(X)$ & Ran space of $X$ \\
$\Ran_{\leq n}(X)$ & Points of $\Ran(X)$ with $\leq n$ elements \\
$\Ran_n(X)$ & Points of $\Ran(X)$ with exactly $n$ elements \\
%
$D_{ij}$ & Boundary divisor $\{z_i = z_j\}$ in $\FM_n$ \\
$D_T$ & Boundary stratum labeled by tree $T$ \\
$\partial \FM_n$ & Total boundary divisor $\bigcup_{i < j} D_{ij}$ \\
$\Omega^k(M)$ & Smooth $k$-forms on manifold $M$ \\
$\Omega^k_{\log}(Y, D)$ & Logarithmic $k$-forms with poles along $D$ \\
%
$\cD'(U)$ & Distributions on open set $U$ \\
$\cD'^k(U)$ & Distributional $k$-currents \\
$\eta_{ij}$ & FM propagator form $\dlog(z_i - z_j)$ \\
$\Res_D$ & Residue along divisor $D$ \\
$\Det(V)$ & Determinant line of vector space $V$ \\
$\orline{M}$ & Orientation line bundle of manifold $M$ \\
$\omega_X$ & Canonical bundle (dualizing sheaf) of $X$ \\
$\omega_{X^n/X}$ & Relative dualizing sheaf \\
\end{longtable}


\section{Chiral structures}
\label{sec:notation-chiral}

\begin{longtable}{p{3cm} p{10cm}}
\caption{Chiral structures} \\
\toprule
\textbf{Symbol} & \textbf{Meaning} \\
\midrule
\endfirsthead
\toprule
\textbf{Symbol} & \textbf{Meaning} \\
\midrule
\endhead
\midrule
\multicolumn{2}{r}{\textit{Continued on next page}} \\
\endfoot
\bottomrule
\endlastfoot

$\cA, \cB, \ldots$ & Chiral algebras \\
$\cC, \cD, \ldots$ & Chiral coalgebras \\
$\mu^{\mathrm{ch}}$ & Chiral bracket/product \\
$\Delta^{\mathrm{ch}}$ & Chiral coproduct \\
$\chirtensor$ & Chiral tensor product \\
$\facttensor$ & Factorization tensor product (same as $\otimes^!$) \\
$j_* j^*$ & Localization away from diagonal \\
$\Delta_!$ & Exceptional direct image along diagonal \\
$\Delta_*$ & Direct image along diagonal \\
$\Chir^{\Eone}$ & Category of $\Eone$-chiral algebras on a curve (associative/nonlocal BD sense) \\
$\Chir^{\Einf}$ & Category of $\Einf$-chiral algebras on a curve (commutative/local; vertex algebras) \\
$\Chir^{\Pinf}$ & Category of $\Pinf$-chiral algebras on a curve \\
%
$\chirLie$ & Chiral Lie operad \\
$\chirAss$ & Chiral associative operad \\
$\chirCom$ & Chiral commutative operad \\
$\chirPois$ & Chiral Poisson operad \\
$\cA^{\text{\textexclamdown}}$ & Koszul dual \emph{coalgebra}: $\cA^{\text{\textexclamdown}} = \Bbar(\cA)$, the reduced bar construction viewed as a conilpotent dg coalgebra.  This is the primary output of the bar construction; no dualization is involved. \\
%
$\cA^!$ & Koszul dual \emph{algebra}: $\cA^! = \Bbar(\cA)^{\vee} = (\cA^{\text{\textexclamdown}})^\vee$, obtained from the Koszul dual coalgebra by graded linear dualization (equivalently $\VD(\Bbar(\cA)) \otimes \omega^{-1}$ in the chiral setting).  Three objects, three symbols: $\Bbar(\cA)$ = bar coalgebra, $\cA^{\text{\textexclamdown}} = \Bbar(\cA)$ = Koszul dual coalgebra, $\cA^! = \Bbar(\cA)^\vee$ = Koszul dual algebra. \\
$H^{\mathrm{ch}}_*(X, \cA)$ & Chiral homology of $\cA$ over $X$ \\
$H^{*,\mathrm{ch}}(X, \cA)$ & Chiral cohomology \\
$\CC^*_{\mathrm{ch}}(\cA)$ & Chiral Hochschild cochain complex \\
$\CC_*^{\mathrm{ch}}(\cA)$ & Chiral Hochschild chain complex \\
%
$\ChirHoch^*(\cA)$ & Chiral Hochschild cohomology \\
$HH_*^{\mathrm{ch}}(\cA)$ & Chiral Hochschild homology \\
$Y(a, z)$ & State-field correspondence/vertex operator \\
$a_{(n)} b$ & $n$-th Borcherds product \\
$a_{(n,m)} b$ & Higher Borcherds product \\
%
$\{a_\lambda b\}$ & $\lambda$-bracket \\
$: ab :$ & Normal ordering \\
$T$ & Translation operator \\
$T(z)$ & Stress-energy tensor (Virasoro field) of conformal weight $2$ \\
$L_n$ & Virasoro modes \\
$c$ & Central charge \\
\end{longtable}


\section{Miscellaneous notation}
\label{sec:notation-misc}

\begin{longtable}{p{3cm} p{10cm}}
\caption{Miscellaneous notation} \\
\toprule
\textbf{Symbol} & \textbf{Meaning} \\
\midrule
\endfirsthead
\toprule
\textbf{Symbol} & \textbf{Meaning} \\
\midrule
\endhead
\midrule
\multicolumn{2}{r}{\textit{Continued on next page}} \\
\endfoot
\bottomrule
\endlastfoot

$k$ & Ground field (typically $\bC$ or $\bQ$) \\
$\bC, \bR, \bQ, \bZ$ & Complex, real, rational numbers; integers \\
$\bA^n$ & Affine $n$-space \\
$\bP^n$ & Projective $n$-space \\
$\bG_m$ & Multiplicative group \\
$\Sigma_n$ & Symmetric group on $n$ letters \\
$\sgn_n$ & Sign representation of $\Sigma_n$ \\
$B_n$ & Braid group on $n$ strands (not to be confused with the unordered configuration space $B_n(X)$) \\
%
$P_n$ & Pure braid group \\
$s$ & Suspension (degree shift by $+1$) \\
$s^{-1}$ & Desuspension (degree shift by $-1$) \\
$V[n]$ & Degree shift (cohomological): $(V[n])^k = V^{k+n}$, so $sV = V[-1]$ and $|sv| = |v| + 1$ \\
$|a|$ & Degree of element $a$ \\
%
$d$ & Differential (generic) \\
$\partial$ & Boundary operator \\
$\delta$ & Overloaded: (i) coboundary or Hochschild differential; (ii) total boundary divisor class on $\overline{\mathcal{M}}_g$ (Mumford's formula $12\lambda = \kappa_1 + \delta$); (iii) Kronecker delta $\delta_{m,n}$ or $\delta_{m+n,0}$.  Meaning is always clear from context and subscripts. \\
$h$ & Conformal weight of a field \\
$\mathsf{h}$ & Coxeter number of a simple Lie algebra (distinguished from conformal weight $h$) \\
%
$h^\vee$ & Dual Coxeter number of a simple Lie algebra \\
$\kappa$ & Level parameter (as subscript, e.g., $\hat{\mathfrak{g}}_\kappa$) \\
$\kappa(\cA)$ & Obstruction coefficient of chiral algebra $\cA$ \\
$\kappa^{ab}$ & Killing form / Casimir tensor on $\mathfrak{g}$ \\
$\kappa_{\mathrm{KS}}$ & Kodaira--Spencer map \\
%
$\kappa_1$ & First Mumford--Morita--Miller class on $\overline{\mathcal{M}}_g$ \\
$\varrho(\mathfrak{g})$ & Anomaly ratio $\sum_{i=1}^r 1/(m_i+1)$ \\
$\sigma$ & Chevalley anti-involution (on $\mathfrak{g}$); also Verdier--Koszul involution (on bar-cobar).  Context always disambiguates: the Chevalley involution acts on generators of~$\fg$, while the Verdier--Koszul involution acts on bar-cobar complexes. \\
$[-, -]$ & Bracket (Lie, graded, etc.) \\
$\{-, -\}$ & Poisson bracket or antibracket \\
%
$\simeq$ & Quasi-isomorphism or equivalence \\
$\cong$ & Isomorphism \\
$\sim$ & Homotopy or weak equivalence \\
$\xrightarrow{\sim}$ & Quasi-isomorphism arrow \\
$\hookrightarrow$ & Inclusion/monomorphism \\
%
$\twoheadrightarrow$ & Surjection/epimorphism \\
$\otimes$ & Tensor product \\
$\boxtimes$ & External tensor product \\
$\wedge$ & Wedge product (exterior) \\
$\circ$ & Composition (operadic or categorical) \\
$\circ_i$ & Partial composition at $i$-th input \\[6pt]
\multicolumn{2}{l}{\emph{Arithmetic shadow tower
  (Chapter~\ref{chap:arithmetic-shadows}):}} \\
$\varepsilon^c_s(\cA)$ & Constrained Epstein zeta:
  $\sum_{\Delta\in\cS}(2\Delta)^{-s}$, summed over scalar
  primary dimensions \\
$\varphi(s)$ & Eisenstein scattering matrix:
  $\Lambda(1{-}s)/\Lambda(s)$ on
  $\mathrm{SL}(2,\bZ)\backslash\mathfrak{H}$ \\
$\Lambda(s)$ & Completed $L$-function:
  $\pi^{-s}\zeta(2s)\Gamma(s)$ \\
$\sigma_{-1}(N)$ & Reciprocal divisor sum:
  $\sum_{d\mid N}1/d$; Dirichlet series
  $\sum\sigma_{-1}(N)N^{-s}=\zeta(s)\zeta(s{+}1)$ \\
$\alpha(\cA)$ & Cusp growth rate of
  $\widehat{Z}^c_\cA$ at the cusp of
  $\cM_{1,1}$ \\
\end{longtable}


\section{Index of key definitions}
\label{sec:index-definitions}

Index of key definitions with page references.

\subsection*{Canonical anchor labels}

\begin{definition}[A-infinity algebra]\label{def:ainf-algebra}
We use the curved/filtered \(A_\infty\) framework fixed in Definition~\ref{def:curved-ainfty}.
\end{definition}

\begin{definition}[Arnold relations]\label{def:arnold-relations}
The Arnold-relation package is the one established in Theorem~\ref{thm:arnold-relations}.
\end{definition}

\begin{definition}[Bar complex (algebraic)]\label{def:bar-complex-algebraic}
The algebraic bar complex is represented by the DG coalgebra model of Theorem~\ref{thm:bar-coalgebra}.
\end{definition}

\begin{definition}[Geometric bar complex]\label{def:geometric-bar-complex}
The geometric bar complex is Definition~\ref{def:geometric-bar}.
\end{definition}

\begin{definition}[Central charge]\label{def:central-charge}
Central charge is the anomaly/curvature parameter appearing in the genus-corrected chiral package; see Example~\ref{ex:virasoro-central-charge-curvature}.
\end{definition}

\begin{definition}[Chiral algebra (BD)]\label{def:chiral-algebra-bd}
We use Beilinson--Drinfeld chiral algebras as implemented in Theorem~\ref{thm:chiral-factorization}.
\end{definition}

\begin{construction}[Chiral bracket]\label{constr:chiral-bracket}
The chiral bracket is encoded by residue/locality operations as in Theorem~\ref{thm:residue-operations}.
\end{construction}

\begin{definition}[Topological chiral homology]\label{def:topological-chiral-homology}
Topological chiral homology is realized in this manuscript via the factorization/bar model of Theorem~\ref{thm:bar-factorization-homology}.
\end{definition}

\begin{definition}[Chiral pseudo-tensor structure]\label{def:chiral-pseudo-tensor}
The pseudo-tensor/factorization tensor structure is fixed by Definition~\ref{def:factorization-algebra-AF}.
\end{definition}

\begin{definition}[Cobar complex]\label{def:cobar-coalgebra}
The cobar complex is defined intrinsically via Verdier duality
(Definition~\ref{def:geom-cobar-intrinsic}); its explicit distributional
model is Definition~\ref{def:geom-cobar}.
\end{definition}

\begin{definition}[Configuration space]\label{def:config-space}
Configuration spaces are as in Definition~\ref{def:config-space-genus-g}.
\end{definition}

\begin{definition}[Classical D-module input]\label{def:dmod-classical}
The \(\mathcal{D}\)-module realization used throughout is the one in Theorem~\ref{thm:chiral-ran}.
\end{definition}

\begin{definition}[E1-chiral algebra]\label{def:e1-chiral-algebra}
An \(\mathbb{E}_1\)-chiral algebra is an associative (but not necessarily commutative)
algebra object in the chiral compound tensor \(\infty\)-category of \(\mathcal{D}_X\)-modules;
see Definition~\ref{def:e1-chiral} for the precise formulation.
\end{definition}

\begin{definition}[E-infinity-chiral algebra]\label{def:einf-chiral-algebra}
An \(\mathbb{E}_\infty\)-chiral algebra is a vertex algebra in the classical sense of Borcherds,
i.e., one satisfying locality (skew-symmetry); see Definition~\ref{def:einf-chiral}.
\end{definition}

\begin{definition}[Factorization algebra]\label{def:factorization-algebra}
Factorization algebras are fixed by Definition~\ref{def:factorization-algebra-AF}.
\end{definition}

\begin{definition}[Fulton--MacPherson compactification]\label{def:fm-compactification}
The compactification package is given by Theorem~\ref{thm:FM}.
\end{definition}

\begin{definition}[Holonomic regime]\label{def:holonomic}
The holonomic regime for chiral/\(\mathcal{D}\)-module objects is used through the Ran-space formalism of Theorem~\ref{thm:chiral-ran}.
\end{definition}

\begin{definition}[Koszul dual algebra]\label{def:koszul-dual-alg}
The intrinsic Koszul dual algebra object is Definition~\ref{def:intrinsic-koszul-dual}.
\end{definition}

\begin{definition}[Koszul dual coalgebra (factorization side)]\label{def:koszul-dual-coalgebra-fact}
The factorization-side dual coalgebra model is Definition~\ref{def:completed-dual}.
\end{definition}

\begin{definition}[Koszul operad]\label{def:koszul-operad-index}
The operadic Koszul mechanism used in this manuscript is represented by Theorem~\ref{thm:operadic-bar}.
\end{definition}

\begin{definition}[Normal ordering]\label{def:normal-ordering}
Normal ordering conventions are those in the explicit OPE package of Definition~\ref{def:beta-gamma-ope-complete}.
\end{definition}

\begin{definition}[Operator product expansion]\label{def:ope}
The canonical affine OPE template is Equation~\eqref{eq:affine-kac-moody-ope}.
\end{definition}

\begin{definition}[Operad as monoid]\label{def:operad-as-monoid}
Operadic monoidal composition for configuration operations is implemented via Theorem~\ref{thm:FM-operad}.
\end{definition}

\begin{definition}[Pro-nilpotent completion]\label{def:pro-nilpotent}
Pro-nilpotent and completed bar behavior is encoded by
Definition~\ref{def:conilpotent-complete}.
\end{definition}

\begin{definition}[Residue map]\label{def:residue-map}
Residue operations are fixed by Theorem~\ref{thm:residue-operations}.
\end{definition}

\begin{definition}[State-field correspondence]\label{def:state-field-correspondence}
State-field correspondence is used in the CFT/chiral bridge formalism of Definition~\ref{def:w-algebra-cft}.
\end{definition}

\begin{definition}[Twisting morphism]\label{def:twisting-morphism-index}
See Definition~\ref{def:twisting-morphism} for the full treatment.
A twisting morphism $\tau\colon C \to A$ is a degree-shifting map
satisfying the MC equation $\partial\tau + \tau \star \tau = 0$ in
the convolution algebra $\mathrm{Hom}(C, A)$
(Definition~\ref{def:convolution-dg-lie}).  The set $\mathrm{Tw}(C, A)$
mediates the bar-cobar adjunction
(Proposition~\ref{prop:universal-twisting-adjunction}).
\end{definition}

\begin{definition}[Verdier duality functor]\label{def:verdier-functor}
Verdier duality is fixed by Theorem~\ref{thm:verdier-bar-cobar}.
\end{definition}

\begin{definition}[Vertex-algebra locality]\label{def:locality}
Locality is encoded by the factorization-locality principle~\ref{princ:factorization-locality}.
\end{definition}

\begin{multicols}{2}
\begin{itemize}[leftmargin=*, itemsep=0pt]
\item $\Ainf$-algebra, Definition~\ref{def:ainf-algebra}
\item Arnold relations, Definition~\ref{def:arnold-relations}
\item Bar complex (algebraic), Definition~\ref{def:bar-complex-algebraic}
\item Bar complex (geometric), Definition~\ref{def:geometric-bar-complex}
\item Central charge, Definition~\ref{def:central-charge}
\item Chiral algebra, Definition~\ref{def:chiral-algebra-bd}
\item Chiral bracket, Construction~\ref{constr:chiral-bracket}
\item Chiral homology, Definition~\ref{def:topological-chiral-homology}
\item Chiral tensor product, Definition~\ref{def:chiral-pseudo-tensor}
\item Cobar complex, Definition~\ref{def:cobar-coalgebra}
\item Configuration space, Definition~\ref{def:config-space}
\item D-module, Definition~\ref{def:dmod-classical}
\item Determinant line, Definition~\ref{def:det-vect}
\item $\Eone$-chiral algebra, Definition~\ref{def:e1-chiral-algebra}
\item $\Einf$-chiral algebra, Definition~\ref{def:einf-chiral-algebra}
\item Factorization algebra, Definition~\ref{def:factorization-algebra}
\item FM compactification, Definition~\ref{def:fm-compactification}
\item Holonomic, Definition~\ref{def:holonomic}
\item Homotopy transfer, Theorem~\ref{thm:htt}
\item $\infty$-operad, Definition~\ref{def:infty-operad}
\item Koszul dual (algebra), Definition~\ref{def:koszul-dual-alg}
\item Koszul pair, Definition~\ref{def:chiral-koszul-pair} (chiral), Definition~\ref{def:koszul-pair-classical} (classical)
\item Koszul dual (coalgebra), Definition~\ref{def:koszul-dual-coalgebra-fact}
\item Koszul operad, Definition~\ref{def:koszul-operad}
\item Logarithmic forms, Definition~\ref{def:log-forms}
\item Maurer--Cartan element, Definition~\ref{def:mc-element}
\item Minimal model, Definition~\ref{def:minimal-model}
\item Normal ordering, Definition~\ref{def:normal-ordering}
\item OPE, Definition~\ref{def:ope}
\item Operad, Definition~\ref{def:operad-as-monoid}
\item Pro-nilpotent, Definition~\ref{def:pro-nilpotent}
\item Ran space, Definition~\ref{def:ran-space}
\item Residue, Definition~\ref{def:residue-map}
\item SDR, Definition~\ref{def:sdr}
\item Spectral sequence, Construction~\ref{constr:ass-graded-ss}
\item State-field correspondence, Definition~\ref{def:state-field-correspondence}
\item Suspension, Definition~\ref{def:suspension}
\item Twisting morphism, Definition~\ref{def:twisting-morphism}
\item Verdier duality, Definition~\ref{def:verdier-functor}
\item Vertex algebra, Definition~\ref{def:locality}
\end{itemize}
\end{multicols}
